% In this file you should put the actual content of the blueprint.
% It will be used both by the web and the print version.
% It should *not* include the \begin{document}
%
% This blueprint is intended to track the Lean development in the repository.
% Each statement below is connected to the corresponding Lean declaration via
%   \lean{...}, \leanok, and \uses{...}.

\begin{theorem}[Smale 1958]
  \label{thm:test_a}
  \lean{test_a}
  \leanok
  \uses{def:test_def}
  There is a homotopy of immersions of $𝕊^2$ into $ℝ^3$ from the inclusion map to
  the antipodal map $a : q \mapsto -q$.
\end{theorem}

\begin{proof}
  \leanok
  \uses{thm:test_b}
  This obviously follows from what we did so far.
\end{proof}

% --------------------------------------------------------------------
% Taylor and Smoothness
% --------------------------------------------------------------------

\section*{TaylorAndSmoothness}
\label{sec:taylor_and_smoothness}

Throughout this section, $E$ is a real Hilbert space, $f : E \to \mathbb{R}$,
$L \in \mathbb{R}_{\ge 0}$ (encoded in Lean as \texttt{NNReal}),
and $s \subseteq E$.

\begin{definition}[$L$-smoothness on a set]
  \label{def:l_smooth_on}
  \lean{Lean4ML.Optimization.LSmoothOn}
  \leanok
  We say that $f$ is $L$-smooth on $s$ if
  \[
    \|\nabla f(x) - \nabla f(y)\| \le L\|x-y\|,\quad \forall x,y \in s.
  \]
  In Lean, this is encoded as \texttt{LipschitzOnWith L (fun x => gradient f x) s}.
\end{definition}

\begin{definition}[Quadratic upper model]
  \label{def:quadratic_upper_bound}
  \lean{Lean4ML.Optimization.QuadraticUpperBound}
  \leanok
  The quadratic upper model on $s$ is the property
  \[
    f(y) \le f(x) + \langle \nabla f(x), y-x \rangle + \frac{L}{2}\|y-x\|^2,
    \quad \forall x,y \in s.
  \]
\end{definition}

\begin{theorem}[First-order Taylor formula along a line]
  \label{thm:taylor_first_order}
  \lean{Lean4ML.Optimization.taylor_first_order}
  \leanok
  \uses{}
  If $f$ is $C^1$, then for all $x,p \in E$,
  \[
    f(x+p) - f(x) = \int_0^1 \langle \nabla f(x + t p), p \rangle\,dt.
  \]
\end{theorem}

\begin{proof}
  \leanok
  \uses{}
  This is proved in Lean by first establishing a vector-valued fundamental theorem
  of calculus along the path $t \mapsto x + t\cdot p$, and then rewriting the
  Fr\'echet derivative in terms of the gradient.
\end{proof}

\begin{lemma}[Pointwise bound on the line-segment integrand]
  \label{lem:lipschitz_gradient_pointwise}
  \lean{Lean4ML.Optimization.lipschitz_gradient_pointwise}
  \leanok
  \uses{def:l_smooth_on}
  Assume $f$ is $L$-smooth on $s$, $x \in s$, and $x + \gamma (y-x) \in s$ with
  $\gamma \ge 0$. Then
  \[
    \left\langle \nabla f\bigl(x+\gamma (y-x)\bigr) - \nabla f(x),\, y-x \right\rangle
    \le L\gamma\|y-x\|^2.
  \]
\end{lemma}

\begin{proof}
  \leanok
  \uses{}
  In Lean, we first convert Lipschitz-on-set smoothness to a norm inequality
  for gradient differences, then apply Cauchy--Schwarz to bound the inner
  product by a product of norms, and finally rewrite the segment displacement
  norm as $\gamma\|y-x\|$ (for $\gamma \ge 0$).
\end{proof}

\begin{theorem}[Quadratic upper bound from smoothness on a convex set]
  \label{thm:l_smooth_quadratic_upper_bound}
  \lean{Lean4ML.Optimization.l_smooth_quadratic_upper_bound}
  \leanok
  \uses{thm:taylor_first_order,lem:lipschitz_gradient_pointwise,def:l_smooth_on,def:quadratic_upper_bound}
  Assume $f$ is $C^1$, $L$-smooth on a convex set $s$, and $x,y \in s$. Then
  \[
    f(y) \le f(x) + \langle \nabla f(x), y-x \rangle + \frac{L}{2}\|y-x\|^2.
  \]
\end{theorem}

\begin{proof}
  \leanok
  \uses{thm:taylor_first_order,lem:lipschitz_gradient_pointwise}
  The Lean proof follows exactly the standard decomposition:
  \begin{enumerate}
    \item apply Theorem~\ref{thm:taylor_first_order} with $p := y-x$ to write
    $f(y)-f(x)$ as an integral of $\langle \nabla f(x+t(y-x)), y-x\rangle$;
    \item subtract the constant term $\langle \nabla f(x), y-x\rangle$ inside the
    integral, rewriting the remainder as an integral of
    $\langle \nabla f(x+t(y-x)) - \nabla f(x), y-x\rangle$;
    \item use convexity of $s$ to ensure $x+t(y-x)\in s$ for all $t\in[0,1]$;
    \item bound the integrand pointwise using Lemma~\ref{lem:lipschitz_gradient_pointwise}
    (with $\gamma=t$) and integrate the bound on $[0,1]$;
    \item evaluate $\int_0^1 Lt\,dt = L/2$ to obtain the desired inequality.
  \end{enumerate}
\end{proof}

% --------------------------------------------------------------------
% Hessian vs smoothness (Lemma 2.3-style statements)
% --------------------------------------------------------------------

\section*{HessianVsSmoothness}
\label{sec:hessian_vs_smoothness}

\begin{definition}[Hessian]
  \label{def:hessian}
  \lean{Lean4ML.Optimization.hessian}
  \leanok
  The Hessian of $f$ at $x$ is the Fr\'echet derivative of the gradient map:
  \[
    \nabla^2 f(x) := D(\nabla f)(x),
  \]
  encoded in Lean as \texttt{fderiv ℝ (fun y => gradient f y) x}.
\end{definition}

\begin{theorem}[Hessian operator-norm bound $\Rightarrow$ global smoothness]
  \label{thm:hessianBoundImpliesLSmoothOnUniv}
  \lean{Lean4ML.Optimization.hessianBoundImpliesLSmoothOnUniv}
  \leanok
  \uses{def:l_smooth_on,def:hessian}
  If the gradient map is differentiable everywhere and its derivative (the Hessian)
  has operator norm uniformly bounded by $L$, then $f$ is $L$-smooth on $\mathbb{R}$-universe
  (i.e.\ on \texttt{Set.univ}).
\end{theorem}

\begin{proof}
  \leanok
  \uses{}
  In Lean, this is a direct application of a Mathlib theorem
  \texttt{lipschitzWith\_of\_nnnorm\_fderiv\_le}, upgraded to the \texttt{Set.univ}
  version via \texttt{lipschitzOnWith\_univ}.
\end{proof}

\begin{theorem}[Global smoothness $\Rightarrow$ Hessian quadratic-form upper bound]
  \label{thm:lSmoothOnUnivImpliesHessianQuadraticFormBound}
  \lean{Lean4ML.Optimization.lSmoothOnUnivImpliesHessianQuadraticFormBound}
  \leanok
  \uses{def:l_smooth_on,def:hessian}
  If $f$ is $L$-smooth on \texttt{Set.univ}, then for all $x,p$,
  \[
    \langle \nabla^2 f(x)p, p\rangle \le L\|p\|^2.
  \]
\end{theorem}

\begin{proof}
  \leanok
  \uses{}
  The Lean proof extracts an operator-norm bound on $\nabla^2 f(x)$ from Lipschitz
  smoothness using the lemma \texttt{opNorm\_fderiv\_le\_of\_lipschitz}, and then
  bounds the quadratic form $\langle (\nabla^2 f(x))p,p\rangle$ by
  $\|\nabla^2 f(x)\|\,\|p\|^2$ via Cauchy--Schwarz and the operator norm inequality.
\end{proof}

% --------------------------------------------------------------------
% Necessary conditions (Theorem 2.4)
% --------------------------------------------------------------------

\section*{NecessaryConditions}
\label{sec:necessary_conditions}

\begin{theorem}[First-order necessary condition]
  \label{thm:local_min_gradient_zero}
  \lean{Lean4ML.Optimization.local_min_gradient_zero}
  \leanok
  \uses{}
  If $f$ is continuously differentiable and $x$ is a local minimizer of $f$, then
  \[
    \nabla f(x) = 0.
  \]
\end{theorem}

\begin{proof}
  \leanok
  \uses{}
  In Lean, we use the Mathlib lemma \texttt{IsLocalMin.hasFDerivAt\_eq\_zero} applied
  to the Fr\'echet derivative of $f$ at $x$, expressed via \texttt{HasGradientAt}.
  Converting the resulting identity in the dual space back to an inner-product
  statement yields $\langle \nabla f(x), \nabla f(x)\rangle = 0$ and hence
  $\nabla f(x)=0$.
\end{proof}

\begin{theorem}[Second-order necessary condition: Hessian is PSD]
  \label{thm:local_min_hessian_psd}
  \lean{Lean4ML.Optimization.local_min_hessian_psd}
  \leanok
  \uses{thm:local_min_gradient_zero,def:hessian,thm:taylor_first_order}
  Suppose $f$ is $C^2$ and $x$ is a local minimizer of $f$. Then for all $p$,
  \[
    0 \le \langle \nabla^2 f(x)p, p\rangle.
  \]
\end{theorem}

\begin{proof}
  \leanok
  \uses{}
  The Lean proof is by contradiction. Assuming $\langle \nabla^2 f(x)p,p\rangle < 0$,
  one uses an exact second-order Taylor expansion with an integral remainder (proved
  in the file \texttt{NecessaryCondition.lean}) to show that for sufficiently small
  $\alpha>0$,
  \[
    f(x+\alpha p) - f(x) < 0,
  \]
  contradicting local minimality. The analytic core is controlling the integral
  remainder using continuity of the Hessian along the line segment.
\end{proof}

% --------------------------------------------------------------------
% Sufficient condition (Theorem 2.5, operational form)
% --------------------------------------------------------------------

\section*{SufficientConditions}
\label{sec:sufficient_conditions}

\begin{definition}[Strict local minimizer]
  \label{def:IsStrictLocalMin}
  \lean{Lean4ML.Optimization.IsStrictLocalMin}
  \leanok
  A strict local minimizer is a point $x$ for which there exists a neighborhood
  $U$ of $x$ such that $f(x) < f(y)$ for all $y\in U$ with $y\neq x$.
\end{definition}

\begin{lemma}[Taylor lower bound from a uniform Hessian lower bound]
  \label{lem:taylor_lower_bound_from_hessian_lb}
  \lean{Lean4ML.Optimization.taylor_lower_bound_from_hessian_lb}
  \leanok
  \uses{def:hessian}
  Assume $f$ is $C^2$ and $\nabla f(x)=0$. If there exist $\varepsilon>0$ and $\rho>0$
  such that for all $y\in B(x,\rho)$ and all $p$,
  \[
    \varepsilon\|p\|^2 \le \langle \nabla^2 f(y)p,p\rangle,
  \]
  then for all $p$ with $\|p\|<\rho$,
  \[
    \frac{\varepsilon}{2}\|p\|^2 \le f(x+p)-f(x).
  \]
\end{lemma}

\begin{proof}
  \leanok
  \uses{}
  This is obtained by integrating the pointwise Hessian lower bound inside the
  exact double-integral remainder form of Taylor's theorem (proved earlier in the
  necessary-condition development), and then evaluating the resulting elementary
  integral on $[0,1]$.
\end{proof}

\begin{theorem}[Second-order sufficient condition (uniform form)]
  \label{thm:strict_local_min_of_gradient_zero_hessian_uniform_lb}
  \lean{Lean4ML.Optimization.strict_local_min_of_gradient_zero_hessian_uniform_lb}
  \leanok
  \uses{def:IsStrictLocalMin,lem:taylor_lower_bound_from_hessian_lb,def:hessian}
  Suppose $f$ is $C^2$ and $\nabla f(x)=0$. If there exist $\varepsilon>0$ and $\rho>0$
  such that for all $y\in B(x,\rho)$ and all $p$,
  \[
    \varepsilon\|p\|^2 \le \langle \nabla^2 f(y)p,p\rangle,
  \]
  then $x$ is a strict local minimizer of $f$.
\end{theorem}

\begin{proof}
  \leanok
  \uses{}
  Take the neighborhood $U := B(x,\rho)$. For any $y\in U$ with $y\neq x$, set
  $p:=y-x$ so that $y=x+p$ and $\|p\|<\rho$. The previous lemma gives a strictly
  positive lower bound on $f(x+p)-f(x)$, hence $f(x)<f(y)$.
\end{proof}

\begin{theorem}[Second-order sufficient condition (positive definite Hessian, finite-dimensional)]
  \label{thm:second_order_sufficient_condition}
  \lean{Lean4ML.Optimization.second_order_sufficient_condition}
  \leanok
  \uses{thm:strict_local_min_of_gradient_zero_hessian_uniform_lb,def:hessian}
  In finite dimensions, if $f$ is $C^2$, $\nabla f(x)=0$, and the Hessian at $x$
  is positive definite in the sense that
  \[
    p\neq 0 \implies 0 < \langle \nabla^2 f(x)p,p\rangle,
  \]
  then $x$ is a strict local minimizer.
\end{theorem}

\begin{proof}
  \leanok
  \uses{}
  In Lean, the proof packages the analytic step ``positive definite at $x$ implies
  a uniform coercivity bound on a small ball around $x$'' (proved using compactness
  of the unit sphere in finite dimensions), and then applies the uniform-form
  sufficient condition.
\end{proof}


% --------------------------------------------------------------------
% Gradient Flow
% --------------------------------------------------------------------

\section*{GradientFlow}
\label{sec:gradient_flow}

\begin{definition}[Gradient flow certificate]
  \label{def:GradientFlow}
  \lean{Lean4ML.Optimization.GradientFlow}
  \leanok
  A \emph{gradient flow} for $f : E \to \mathbb{R}$ with initial point $x_0 \in E$
  on a time interval $[0,T]$ is a trajectory $x : \mathbb{R}\to E$ such that
  $x(0)=x_0$ and $\dot x(t) = -\nabla f(x(t))$ for all $t\in[0,T]$.
\end{definition}

\begin{theorem}[Gradient flow existence on a finite time horizon]
  \label{thm:gradientFlow_existence}
  \lean{Lean4ML.Optimization.gradientFlow_existence}
  \leanok
  \uses{def:GradientFlow,def:l_smooth_on}
  If $f$ is $C^1$, globally $L$-smooth, $T>0$, and $x_0\in E$, then there exists
  a gradient flow solution on $[0,T]$ starting at $x_0$.
\end{theorem}

\begin{proof}
  \leanok
  \uses{}
  The Lean proof proceeds in three steps.
  \begin{enumerate}
    \item (\emph{Uniform local existence}) Using Picard--Lindel\"of for the ODE
    $\dot x(t) = -\nabla f(x(t))$, global $L$-smoothness gives a uniform step size
    $\delta = \frac{1}{4L+4}$ such that for every starting point $x_0$ there is a
    solution on $[0,\delta]$.
    \item (\emph{Iteration / extension}) By concatenating these local solutions,
    we obtain a solution on $[0,n\delta]$ for any $n\in\mathbb{N}$.
    \item (\emph{Restriction}) Taking $n$ large enough so that $T \le n\delta$ and
    restricting the trajectory yields a solution on $[0,T]$.
  \end{enumerate}
\end{proof}

% --------------------------------------------------------------------
% Gradient Flow (unbounded time)
% --------------------------------------------------------------------

\begin{definition}[Global gradient flow]
  \label{def:GradientFlowGlobal}
  \lean{Lean4ML.Optimization.GradientFlowGlobal}
  \leanok
  A \emph{global} gradient flow for $f : E \to \mathbb{R}$ with initial point $x_0$
  is a trajectory $x : \mathbb{R}\to E$ such that $x(0)=x_0$ and
  $\dot x(t) = -\nabla f(x(t))$ for all $t \ge 0$.
\end{definition}

\begin{theorem}[Global gradient flow existence for all $t \ge 0$]
  \label{thm:gradientFlowGlobal_existence}
  \lean{Lean4ML.Optimization.gradientFlowGlobal_existence}
  \leanok
  \uses{def:GradientFlowGlobal,thm:gradientFlow_existence}
  If $f$ is $C^1$ and globally $L$-smooth, then for every $x_0\in E$ there exists
  a global gradient flow solution defined for all $t \ge 0$.
\end{theorem}

\begin{proof}
  \leanok
  \uses{}
  For each $n\in\mathbb{N}$, apply the finite-horizon existence theorem to obtain
  a solution on $[0,n+1]$. Uniqueness of solutions on overlapping intervals implies
  these flows are consistent under restriction. Define the global trajectory by
  \[
    x(t) := x_{\lceil t\rceil}(t),
  \]
  i.e.\ evaluate the $(\lceil t\rceil)$-th finite-horizon flow at time $t$.
  A locality (\emph{eventually equal}) argument shows this global definition agrees
  with one fixed finite-horizon solution in a neighborhood of each $t\ge 0$, so
  differentiability and the ODE identity transfer pointwise, yielding a global
  gradient flow.
\end{proof}
